% Source-faithful mathematical blueprint of Hartshorne, Algebraic Geometry (1977), Chapter V.
% Authorial exposition, examples, and exercises are omitted; definitions, statements, and proof
% arguments are retained. The order and numbering below are those of Chapter V.

\section{Geometry on a Surface}
\label{ha-v1}
\source{hartshorne-algebraic-geometry:V.1}

\begin{convention}[Surfaces, curves, and numerical equivalence]
\source{hartshorne-algebraic-geometry:V.1}
Throughout V, a surface is a nonsingular projective surface over an algebraically
closed field, and a curve is an effective divisor unless a different meaning is
specified.  A point is closed, divisors are Cartier divisors, linear equivalence
is denoted by \(\sim\), and \(\operatorname{Num}X\) is the group of divisors
modulo numerical equivalence.
\end{convention}

\begin{definition}[Local intersection, self-intersection, and canonical divisor]
\source{hartshorne-algebraic-geometry:V.1}
For curves $C,D$ with no common component, define the local intersection
multiplicity at $P$ by
\[
(C.D)_P=\length_{\mathcal O_{X,P}}\mathcal O_{X,P}/(f,g),
\]
for local equations $f,g$.  The self-intersection of a nonsingular curve is
$C^2=\deg(\mathcal O_X(C)|_C)$, equivalently the degree of its normal bundle.
The canonical divisor $K_X$ is represented by
$\omega_X=\bigwedge^2\Omega_{X/k}$.
\end{definition}

\begin{theorem}[Intersection pairing]
\label{ha-v1-1}\dcref{V.1.1}\source{hartshorne-algebraic-geometry:V.1.1}
There is a unique pairing
\[
 \operatorname{Div}X\times\operatorname{Div}X\longrightarrow\mathbb Z,
 \qquad (C,D)\longmapsto C.D,
\]
such that: (i) if nonsingular curves meet transversally, \(C.D\) is the number
of their intersection points; (ii) \(C.D=D.C\); (iii) it is additive in each
variable; and (iv) it depends only on the linear-equivalence classes of the two
divisors.
\uses{ha-v1-2,ha-v1-3}
\end{theorem}
\begin{proof}
For a very ample divisor \(A\), Bertini gives a member of \(|A|\) that is
nonsingular and transverse to any prescribed finite collection of curves
(Lemma~\ref{ha-v1-2}).  Write arbitrary divisors as differences
\(C\sim C'-E'\), \(D\sim D'-F'\), with all four terms nonsingular and chosen
transversally.  Set
\[
 C.D=C'.D'-C'.F'-E'.D'+E'.F'.
\]
Replacing one member by another in the same very ample system changes the four
intersection counts by equal amounts, by Lemma~\ref{ha-v1-3}; hence the signed
count is well defined.  It is visibly additive and symmetric, and agrees with
the prescribed transverse count.  Every divisor is a difference of such curves,
so uniqueness follows.
\end{proof}

\begin{lemma}[Bertini choice]
\label{ha-v1-2}\dcref{V.1.2}\source{hartshorne-algebraic-geometry:V.1.2}
If \(C_1,\ldots,C_r\) are irreducible curves and \(A\) is very ample, a general
member of \(|A|\) is irreducible, nonsingular, and transverse to every \(C_i\).

\end{lemma}
\begin{proof}
Embed \(X\) by \(|A|\).  Bertini gives a nonsingular irreducible general
hyperplane section.  For each \(C_i\), the hyperplanes tangent to the embedded
curve or containing a non-transverse tangent direction form a proper closed
subset of the dual projective space.  Avoiding their finite union gives the
assertion.
\end{proof}

\begin{lemma}[Intersection with a smooth curve]
\label{ha-v1-3}\dcref{V.1.3}\source{hartshorne-algebraic-geometry:V.1.3}
If \(C\) is nonsingular and \(D\) is transverse to it, then
\[
 C.D=\deg_C\bigl(\mathcal L(D)\otimes\mathcal O_C\bigr).
\]

\end{lemma}
\begin{proof}
The defining equation of \(D\) restricts to a section of
\(\mathcal L(D)|_C\), whose zero scheme is \(C\cap D\).  Transversality makes
the zero scheme reduced, and the exact sequence
\[
0\to\mathcal L(-D)|_C\to\mathcal O_C\to\mathcal O_{C\cap D}\to0
\]
has Euler characteristic equal to the length of this zero scheme.  On a smooth
curve, that length is the degree of \(\mathcal L(D)|_C\).
\end{proof}

\begin{proposition}[Local intersection formula]
\label{ha-v1-4}\dcref{V.1.4}\source{hartshorne-algebraic-geometry:V.1.4}
If curves \(C,D\) have no common component, then
\[
 C.D=\sum_{P\in C\cap D}(C.D)_P,qquad
 (C.D)_P=\length_{\mathcal O_{X,P}}\mathcal O_{X,P}/(f,g),
\]
where \(f,g\) are local equations at \(P\).
\uses{ha-v1-1}
\end{proposition}
\begin{proof}
Choose a nonsingular curve linearly equivalent to one divisor and transverse
to the other, and use Theorem~\ref{ha-v1-1}.  The displayed exact sequence,
now interpreted on the finite scheme \(C\cap D\), identifies its Euler
characteristic with the sum of the local lengths.  Additivity and linear
equivalence then give the formula for arbitrary \(C,D\).
\end{proof}

\begin{proposition}[Adjunction formula]
\label{ha-v1-5}\dcref{V.1.5}\source{hartshorne-algebraic-geometry:V.1.5}
For a nonsingular curve \(C\) of genus \(g\) on \(X\),
\[
 2g-2=C.(C+K_X).
\]

\end{proposition}
\begin{proof}
The conormal sequence and the determinant of the tangent sequence give
\(\omega_C\simeq(\omega_X\otimes\mathcal O_X(C))|_C\).  Taking degrees and
using \(\deg\omega_C=2g-2\) yields the formula.
\end{proof}

\begin{theorem}[Riemann--Roch for surfaces]
\label{ha-v1-6}\dcref{V.1.6}\source{hartshorne-algebraic-geometry:V.1.6}
Put \(l(D)=h^0(X,\mathcal L(D))\), \(s(D)=h^1(X,\mathcal L(D))\), and
\(p_a=\chi(\mathcal O_X)-1\).  Then
\[
 l(D)-s(D)+l(K_X-D)=\frac12D.(D-K_X)+1+p_a.
\]
\uses{ha-v1-5}
\end{theorem}
\begin{proof}
Serre duality gives \(h^2(\mathcal L(D))=l(K_X-D)\), so it suffices to
compute \(\chi(\mathcal L(D))\).  Write \(D=C-E\), with \(C,E\) effective
and without common components, and use
\(0\to\mathcal L(-E)\to\mathcal O_X\to\mathcal O_E\to0\) and its twist by
\(\mathcal L(C)\).  Additivity of Euler characteristics reduces the calculation
to \(\chi(\mathcal O_C)\) and \(\chi(\mathcal O_E)\).  Curve Riemann--Roch,
adjunction, and Theorem~\ref{ha-v1-5} give
\[
 \chi(\mathcal L(D))=1+p_a+\tfrac12D.(D-K_X),
\]
which is the assertion.
\end{proof}

\begin{lemma}[Vanishing in large positive degree]
\label{ha-v1-7}\dcref{V.1.7}\source{hartshorne-algebraic-geometry:V.1.7}
If \(H\) is ample, there is an integer \(n_0\) such that
\(H^2(X,\mathcal L(D))=0\) whenever \(D.H>n_0\).

\end{lemma}
\begin{proof}
By duality this is equivalent to
\(H^0(X,\mathcal L(K_X-D))=0\).  Choose \(n_0\) larger than the \(H\)-degree
of every component of an effective representative of \(K_X-D\) that can
occur in a fixed bounded family.  If a section existed, \(K_X-D\) would be
effective and would have bounded positive \(H\)-degree, contradicting
\(D.H>n_0\).  The boundedness follows from ampleness and the finite-dimensional
spaces of divisors of bounded degree.
\end{proof}

\begin{corollary}[Effectivity criterion]
\label{ha-v1-8}\dcref{V.1.8}\source{hartshorne-algebraic-geometry:V.1.8}
If \(D.H>0\) and \(D^2>0\), then \(nD\) is linearly equivalent to an effective
divisor for all sufficiently large \(n\).
\uses{ha-v1-6,ha-v1-7}
\end{corollary}
\begin{proof}
Apply Theorem~\ref{ha-v1-6} to \(nD\).  Lemma~\ref{ha-v1-7} makes the second
cohomology term vanish for large \(n\), while the quadratic term
\(\frac12n^2D^2\) dominates the linear term and the fixed constant.  Thus
\(h^0(nD)>0\), so \(|nD|\ne\varnothing\).
\end{proof}

\begin{theorem}[Hodge index theorem]
\label{ha-v1-9}\dcref{V.1.9}\source{hartshorne-algebraic-geometry:V.1.9}
If \(H\) is ample and \(D\not\equiv0\) with \(D.H=0\), then \(D^2<0\).
Equivalently, the intersection form on \(\operatorname{Num}X_\mathbb R\) has
signature \((1,\rho-1)\).
\uses{ha-v1-8}
\end{theorem}
\begin{proof}
For small rational \(\epsilon>0\), apply Corollary~\ref{ha-v1-8} to
\(D+\epsilon H\) and to suitable integral multiples.  If \(D^2>0\), the
resulting effective divisors contradict positivity against the ample class
after subtracting their \(H\)-components.  If \(D^2=0\), perturbation and the
same argument produce a nonzero class orthogonal to \(H\) with negative square,
and equality forces numerical proportionality to \(H\), contrary to
\(D.H=0\).  Hence the orthogonal complement of \(H\) is negative definite.
\end{proof}

\begin{theorem}[Nakai--Moishezon criterion]
\label{ha-v1-10}\dcref{V.1.10}\source{hartshorne-algebraic-geometry:V.1.10}
A divisor \(D\) on a nonsingular projective surface is ample if and only if
\(D^2>0\) and \(D.C>0\) for every irreducible curve \(C\).

\end{theorem}
\begin{proof}
Necessity follows from positivity of an ample restriction and of its top
self-intersection.  Conversely replace \(D\) by a linearly equivalent effective
divisor and use the Hodge index theorem to show that \(D\) has positive degree
on every component of its reduced support.  The restriction of a sufficiently
large multiple to that support is ample; the exact sequence
\(0\to\mathcal O_X(nD-D_{\mathrm{red}})\to\mathcal O_X(nD)\to
\mathcal O_{D_{\mathrm{red}}}(nD)\to0\), together with Serre vanishing, makes
\(|nD|\) separate points and tangent vectors off the support as well.  The
associated morphism has finite fibers, hence is finite by Stein factorization;
its pullback of an ample divisor is a positive multiple of \(D\), so \(D\) is
ample.
\end{proof}

\section{Ruled Surfaces}
\label{ha-v2}\source{hartshorne-algebraic-geometry:V.2}

\begin{definition}[Geometrically ruled surface]
\source{hartshorne-algebraic-geometry:V.2}
A geometrically ruled surface is a surface \(\pi:X\to C\) with a section and
all fibers \(f\simeq\mathbb P^1\).  For a divisor \(D\), write
\(\mathcal L(D)=\mathcal O_X(D)\), and let \(g=g(C)\).
\end{definition}

\begin{notation}[Normalized ruled notation]
\source{hartshorne-algebraic-geometry:V.2}
For a normalized presentation $X=\mathbb P_C(\mathcal E)$, put
$e=-\deg\mathcal E$ and choose the section $C_0$ with
$\mathcal O_X(C_0)\simeq\mathcal O_X(1)$.  Let $f$ be a fiber.  Then every
numerical divisor class is $aC_0+bf$, with $C_0^2=-e$, $C_0.f=1$, and
$f^2=0$.
\end{notation}

\begin{lemma}[Direct images on fibers]
\label{ha-v2-1}\dcref{V.2.1}\source{hartshorne-algebraic-geometry:V.2.1}
If \(D.f=n\ge0\), then \(\pi_*\mathcal L(D)\) is locally free of rank \(n+1\),
\(R^1\pi_*\mathcal L(D)=0\), and \(\pi_*\mathcal O_X=\mathcal O_C\).

\end{lemma}
\begin{proof}
The restriction to each fiber is \(\mathcal O_{\mathbb P^1}(n)\), whose
\(H^0\) has dimension \(n+1\) and whose \(H^1\) vanishes.  Grauert's theorem
then gives local freeness and base change; applying the same argument to
\(\mathcal O_X\) gives the last assertion.
\end{proof}

\begin{proposition}[Projective-bundle description]
\label{ha-v2-2}\dcref{V.2.2}\source{hartshorne-algebraic-geometry:V.2.2}
Every ruled surface is \(\mathbb P_C(\mathcal E)\) for a locally free sheaf
\(\mathcal E\) of rank two.  Conversely \(\mathbb P_C(\mathcal E)\) is ruled,
and \(\mathbb P(\mathcal E)\simeq\mathbb P(\mathcal E')\) over \(C\) exactly
when \(\mathcal E'\simeq\mathcal E\otimes\mathcal L\) for an invertible sheaf
\(\mathcal L\).

\end{proposition}
\begin{proof}
A section gives a quotient \(\mathcal E\twoheadrightarrow\mathcal L\) on
each trivializing open set.  The resulting morphism to \(\mathbb P(\mathcal E)\)
is an isomorphism on every fiber and hence an isomorphism globally.  Conversely
local triviality of \(\mathcal E\) gives fibers \(\mathbb P^1\) and sections.
The relative \(\mathcal O(1)\) is determined up to tensoring by a line bundle,
which proves the isomorphism criterion.
\end{proof}

\begin{proposition}[Picard and numerical groups]
\label{ha-v2-3}\dcref{V.2.3}\source{hartshorne-algebraic-geometry:V.2.3}
For \(X=\mathbb P_C(\mathcal E)\),
\(\operatorname{Pic}X\simeq\mathbb Z[\mathcal O_X(1)]\oplus\pi^*\operatorname{Pic}C\).
After choosing a section \(C_0\), \(\operatorname{Num}X\simeq\mathbb Z C_0\oplus\mathbb Z f\),
with \(C_0.f=1\) and \(f^2=0\).

\end{proposition}
\begin{proof}
Subtracting a suitable multiple of a fiber makes a divisor have degree zero
on every fiber; its restriction is then pulled back from \(C\).  This gives
surjectivity, while restriction to a fiber proves uniqueness.  Numerical
equivalence on pullbacks is degree on \(C\), yielding the stated basis and
intersection numbers.
\end{proof}

\begin{lemma}[Cohomology of nonnegative fiber degree]
\label{ha-v2-4}\dcref{V.2.4}\source{hartshorne-algebraic-geometry:V.2.4}
If \(D.f\ge0\), then \(R^i\pi_*\mathcal L(D)=0\) for \(i>0\), and
\[
 H^i(X,\mathcal L(D))\simeq H^i(C,\pi_*\mathcal L(D)).
\]

\end{lemma}
\begin{proof}
The fiber restriction has no higher cohomology, so Grauert gives the vanishing
of higher direct images.  The Leray spectral sequence therefore degenerates.
\end{proof}

\begin{corollary}[Invariants of a ruled surface]
\label{ha-v2-5}\dcref{V.2.5}\source{hartshorne-algebraic-geometry:V.2.5}
For a ruled surface over a curve of genus \(g\),
\[
 p_a(X)=-g,\qquad p_g(X)=0,\qquad q(X)=g.
\]
\uses{ha-v2-4}
\end{corollary}
\begin{proof}
Use Lemma~\ref{ha-v2-4} with \(D=0\) and \(\pi_*\mathcal O_X=\mathcal O_C\),
then the cohomology of \(\mathcal O_C\) gives \(h^0=1,h^1=g,h^2=0\).
\end{proof}

\begin{proposition}[Sections and quotient bundles]
\label{ha-v2-6}\dcref{V.2.6}\source{hartshorne-algebraic-geometry:V.2.6}
A section \(D\subset\mathbb P_C(\mathcal E)\) is equivalent to a surjection
\(\mathcal E\twoheadrightarrow\mathcal L\) onto an invertible sheaf.  If
\(0\to\mathcal N\to\mathcal E\to\mathcal L\to0\), then
\[
 \mathcal N\simeq\pi_*(\mathcal O_X(1)\otimes\mathcal O_X(-D)),qquad
 \pi_*\mathcal N\simeq\mathcal O_X(1)\otimes\mathcal O_X(-D)
\]
under the natural identification on the base.
\uses{ha-v2-4}
\end{proposition}
\begin{proof}
The quotient defines the section by the universal property of \(\mathbb P(\mathcal E)\).
The exact sequence
\(0\to\mathcal O_X(1)\otimes\mathcal L^{-1}\to\mathcal O_X(1)\to
\mathcal O_X(1)|_D\to0\), pushed down using Lemma~\ref{ha-v2-4}, identifies
the kernel and its pushforward as displayed.
\end{proof}

\begin{corollary}[Extensions on curves]
\label{ha-v2-7}\dcref{V.2.7}\source{hartshorne-algebraic-geometry:V.2.7}
Every locally free sheaf of rank two on a smooth curve is an extension of two
invertible sheaves.

\end{corollary}
\begin{proof}
Choose a line bundle quotient of the generic fiber and extend it after twisting
by a sufficiently positive line bundle.  The kernel is torsion-free of rank one,
hence invertible on a nonsingular curve, and gives the extension.
\end{proof}

\begin{proposition}[Normalization and the invariant]
\label{ha-v2-8}\dcref{V.2.8}\source{hartshorne-algebraic-geometry:V.2.8}
Every ruled surface has a normalized presentation \(X=\mathbb P(\mathcal E)\),
where \(H^0(\mathcal E)\ne0\) but \(H^0(\mathcal E\otimes\mathcal L)=0\) for every
line bundle \(\mathcal L\) of negative degree.  The integer
\(e=-\deg\mathcal E\) is presentation-independent, and there is a section
\(C_0\) with \(\mathcal O_X(C_0)\simeq\mathcal O_X(1)\).

\end{proposition}
\begin{proof}
Twist \(\mathcal E\) by a line bundle of minimal degree for which it has a
nonzero section.  Minimality gives the vanishing for negative twists.  A nonzero
section yields a quotient after removing its zero divisor and hence the section
\(C_0\).  A further twist changes \(\deg\mathcal E\) by twice the degree of
the twist and changes the chosen section correspondingly, so \(e\) is intrinsic.
\end{proof}

\begin{proposition}[Classes of sections]
\label{ha-v2-9}\dcref{V.2.9}\source{hartshorne-algebraic-geometry:V.2.9}
If a section \(D\) corresponds to \(\mathcal E\twoheadrightarrow\mathcal L\), then
\[
 \deg\mathcal L=C_0.D,\qquad D\sim C_0+(\deg\mathcal L-e)f,
 \qquad C_0^2=-e.
\]
\uses{ha-v2-6}
\end{proposition}
\begin{proof}
The first equality is restriction to the section.  Apply Proposition~\ref{ha-v2-6}
to the kernel and compare degrees in the resulting exact sequence; the divisor
class differs from \(C_0\) by the indicated pullback.  Intersecting with \(C_0\)
then gives \(C_0^2=-e\).
\end{proof}

\begin{lemma}[Canonical divisor of a ruled surface]
\label{ha-v2-10}\dcref{V.2.10}\source{hartshorne-algebraic-geometry:V.2.10}
For the normalized notation, with \(K_C\) also denoting a divisor pulled back
from \(C\),
\[
 K_X\sim -2C_0+(K_C+ef).
\]

\end{lemma}
\begin{proof}
Write \(K_X\sim aC_0+bf\).  Adjunction on a fiber gives
\(-2=K_X\mathbin{.}f=-2\), so \(a=-2\).  Adjunction on \(C_0\), whose genus
is \(g(C)\), gives \(2g-2=C_0\mathbin{.}(C_0+K_X)=-e+b\), hence
\(b=2g-2+e\), which is the displayed
linear equivalence.
\end{proof}

\begin{corollary}[Canonical self-intersection]
\label{ha-v2-11}\dcref{V.2.11}\source{hartshorne-algebraic-geometry:V.2.11}
\[
 K_X\equiv-2C_0+(2g-2-e)f,qquad K_X^2=8(1-g).
\]
\uses{ha-v2-10}
\end{corollary}
\begin{proof}
Use Lemma~\ref{ha-v2-10}, \(C_0^2=-e\), \(C_0.f=1\), and \(f^2=0\), then
expand the square.
\end{proof}

\begin{theorem}[Normalized bundles]
\label{ha-v2-12}\dcref{V.2.12}\source{hartshorne-algebraic-geometry:V.2.12}
For a normalized rank-two bundle on a genus-\(g\) curve, if it is decomposable
then \(\mathcal E\simeq\mathcal O_C\oplus\mathcal L\) with \(\deg\mathcal L\le0\),
so \(e\ge0\), and every \(e\ge0\) occurs.  If it is indecomposable, then
\[
 -2g\le e\le2g-2.
\]
\uses{ha-v2-7}
\end{theorem}
\begin{proof}
Normalization orders the two summands and gives the decomposable assertion.
For an indecomposable bundle use Proposition~\ref{ha-v2-7} in a maximal line
subbundle extension \(0\to\mathcal O_C\to\mathcal E\to\mathcal L\to0\).
The extension class is nonzero in \(H^1(\mathcal L^{-1})\), so Serre duality
and Riemann--Roch force both \(h^0(K_C\otimes\mathcal L)\) and
\(h^0(\mathcal L^{-1})\) to be nonzero only within the stated degree range.
\end{proof}

\begin{corollary}[Rational ruled surfaces]
\label{ha-v2-13}\dcref{V.2.13}\source{hartshorne-algebraic-geometry:V.2.13}
Over \(C=\mathbb P^1\), every ruled surface is isomorphic to exactly one
\(X_e=\mathbb P(\mathcal O\oplus\mathcal O(-e))\), with \(e\ge0\).
\uses{ha-v2-8}
\end{corollary}
\begin{proof}
Every bundle splits on \(\mathbb P^1\), and tensoring makes the larger summand
degree zero.  Proposition~\ref{ha-v2-8} makes the difference nonnegative and
unique.
\end{proof}

\begin{corollary}[Splitting on the projective line]
\label{ha-v2-14}\dcref{V.2.14}\source{hartshorne-algebraic-geometry:V.2.14}
Every locally free rank-two sheaf on \(\mathbb P^1\) is a direct sum of line
bundles.
\uses{ha-v2-7}
\end{corollary}
\begin{proof}
Apply the extension statement of Corollary~\ref{ha-v2-7}; the classification
of line bundles and the vanishing of the relevant \(H^1\) after ordering the
degrees show that the extension splits.
\end{proof}

\begin{theorem}[Indecomposable elliptic ruled surfaces]
\label{ha-v2-15}\dcref{V.2.15}\source{hartshorne-algebraic-geometry:V.2.15}
For an elliptic curve, an indecomposable normalized rank-two bundle has
\(e=0\) or \(e=-1\), and there is exactly one ruled-surface isomorphism class
for each value.

\end{theorem}
\begin{proof}
The extension \(0\to\mathcal O_C\to\mathcal E\to\mathcal L\to0\) and
\(h^1(\mathcal O_C)=1\) show that the nontrivial degree-zero extension is unique.
For degree one, tensoring by a degree-zero line bundle identifies all nontrivial
extensions with the unique normalized one.  If \(e\le-2\), the cohomology
diagram induced by multiplication by a nonzero section of a degree-one line
bundle makes an injective map from a two-dimensional space into a one-dimensional
space, a contradiction.  Twisting gives the asserted isomorphism classes.
\end{proof}

\begin{corollary}[Atiyah's elliptic classification]
\label{ha-v2-16}\dcref{V.2.16}\source{hartshorne-algebraic-geometry:V.2.16}
Indecomposable rank-two bundles of each fixed degree on an elliptic curve are
parametrized by the points of the curve (after the natural determinant/twist
identifications).
\uses{ha-v2-15}
\end{corollary}
\begin{proof}
Twist by a line bundle to reduce the degree to zero or one.  Theorem~\ref{ha-v2-15}
classifies the normalized extensions, and degree-zero line bundles are
parametrized by \(\operatorname{Pic}^0(C)\simeq C\).  Untwisting restores the
classification in every degree.
\end{proof}

\begin{theorem}[Linear systems on \(X_e\)]
\label{ha-v2-17}\dcref{V.2.17}\source{hartshorne-algebraic-geometry:V.2.17}
Let \(X_e=\mathbb P(\mathcal O\oplus\mathcal O(-e))\).  A section
\(D\sim C_0+nf\) exists exactly when \(n=0\) or \(n\ge e\); for \(n=e\) there
is a disjoint section \(C_1\sim C_0+ef\) with \(C_1^2=e\).  Moreover
\(|C_0+nf|\) is base-point-free iff \(n\ge e\), and very ample iff \(n>e\).
\uses{ha-v2-6}
\end{theorem}
\begin{proof}
By Proposition~\ref{ha-v2-6}, a section in this class is a quotient
\(\mathcal O\oplus\mathcal O(-e)\twoheadrightarrow\mathcal O(n-e)\).
Nonzero components exist precisely for \(n=0\) or \(n\ge e\), and for \(n=e\)
the second component gives the disjoint section.  Intersecting with \(C_0\)
shows base-point-freeness exactly for \(n\ge e\).  When \(n>e\), the two
summands separate points and tangent vectors on a fiber, on the minimal section,
and between fibers; the vanishing of
\(H^1(\mathcal O(n-1))\) and \(H^1(\mathcal O(n-e-1))\) supplies the required
lifting.  For \(n=e\), the disjoint section is contracted, so very ampleness fails.
\end{proof}

\begin{corollary}[Ample and very ample classes on \(X_e\)]
\label{ha-v2-18}\dcref{V.2.18}\source{hartshorne-algebraic-geometry:V.2.18}
For \(D=aC_0+bf\) on \(X_e\), \(D\) is ample iff it is very ample iff
\(a>0\) and \(b>ae\).  A general member of such a system is irreducible and
nonsingular.
\uses{ha-v2-17}
\end{corollary}
\begin{proof}
Write \(D=a(C_0+nf)+r f\) and apply Theorem~\ref{ha-v2-17}; the inequalities
are exactly those for a sufficiently positive multiple of the generating
systems.  Necessity follows by intersecting with \(f\) and \(C_0\).  Bertini
gives the final assertion.
\end{proof}

\begin{corollary}[Rational scrolls]
\label{ha-v2-19}\dcref{V.2.19}\source{hartshorne-algebraic-geometry:V.2.19}
For \(n\ge e\), \(|C_0+nf|\) embeds \(X_e\) as a rational scroll of degree
\(2n-e\) in \(\mathbb P^{2n-e+1}\).
\uses{ha-v2-17,ha-v2-1}
\end{corollary}
\begin{proof}
Theorem~\ref{ha-v2-17} gives the embedding.  Intersection theory gives
\((C_0+nf)^2=2n-e\).  Lemma~\ref{ha-v2-1} and cohomology on \(\mathbb P^1\) give
\(h^0(C_0+nf)=2n-e+2\), hence the stated projective dimension and degree.
\end{proof}

\begin{proposition}[Curves and ampleness for \(e\ge0\)]
\label{ha-v2-20}\dcref{V.2.20}\source{hartshorne-algebraic-geometry:V.2.20}
If \(e\ge0\) and an irreducible curve \(Y\) is neither \(C_0\) nor a fiber,
then \(Y\sim aC_0+bf\) with \(a>0\) and \(b\ge ae\).  A divisor
\(D=aC_0+bf\) is ample iff \(a>0\) and \(b>ae\).

\end{proposition}
\begin{proof}
Intersect with \(f\) to get \(a>0\), and with \(C_0\) to get
\(b-ae\ge0\).  Nakai--Moishezon applied to \(f,C_0\), and these inequalities
for every other irreducible curve, gives the criterion.
\end{proof}

\begin{proposition}[Curves and ampleness for \(e<0\)]
\label{ha-v2-21}\dcref{V.2.21}\source{hartshorne-algebraic-geometry:V.2.21}
Assume \(e<0\) and either characteristic zero or \(g\le1\).  If
\(Y\sim aC_0+bf\) is irreducible and is neither a section nor a fiber, then
\(a>0\) and \(b\ge ae/2\).  A divisor \(D=aC_0+bf\) is ample iff
\(a>0\) and \(b>ae/2\).
\uses{ha-v2-9}
\end{proposition}
\begin{proof}
Normalize \(Y\).  Hurwitz gives \(g(\widetilde Y)\ge g(C)\), while adjunction and
\(C_0^2=-e\) give
\((a-1)(b-ae/2)\ge0\).  The case \(a=1\) is handled by the section
description of Proposition~\ref{ha-v2-9}.  These inequalities, together with
the intersections with a fiber and with the minimal section, imply the stated
positivity; Nakai--Moishezon proves sufficiency and necessity.
\end{proof}

\section{Monoidal Transformations}
\label{ha-v3}\source{hartshorne-algebraic-geometry:V.3}

\begin{construction}[Monoidal transformation and multiplicity]
\source{hartshorne-algebraic-geometry:V.3}
The monoidal transformation at $P\in X$ is the blow-up
$\pi:\widetilde X\to X$.  Its inverse image $E=\pi^{-1}(P)$ is the
exceptional curve.  If $C$ has local equation $f$ at $P$, its multiplicity
$\mu_P(C)$ is the largest integer $r$ such that
$f\in\mathfrak m_P^r$.
\end{construction}

\begin{proposition}[The blow-up of a point]
\label{ha-v3-1}\dcref{V.3.1}\source{hartshorne-algebraic-geometry:V.3.1}
If \(\pi:\widetilde X\to X\) is the blow-up of a nonsingular point \(P\) on a
nonsingular projective surface, then \(\widetilde X\) is nonsingular and
projective, the exceptional curve \(E=\pi^{-1}(P)\) is \(\mathbb P^1\), and
\(E^2=-1\).

\end{proposition}
\begin{proof}
In regular parameters \((x,y)\), the blow-up is covered by
\(\operatorname{Spec}k[x,y/x]\) and \(\operatorname{Spec}k[x/y,y]\), so it is
nonsingular and projective.  The exceptional divisor is \(\operatorname{Proj}
\operatorname{gr}_{(x,y)}\mathcal O_{X,P}\simeq\mathbb P^1\).  Its normal bundle
is \(\mathcal O_{\mathbb P^1}(-1)\), hence its self-intersection is \(-1\).
\end{proof}

\begin{proposition}[Picard group and intersections after blowing up]
\label{ha-v3-2}\dcref{V.3.2}\source{hartshorne-algebraic-geometry:V.3.2}
\[
 \operatorname{Pic}\widetilde X\simeq\pi^*\operatorname{Pic}X\oplus\mathbb Z[E],
 \quad (\pi^*C).(\pi^*D)=C.D,\quad \pi^*C.E=0,\quad E^2=-1,
\]
and \(C.\pi_*D=(\pi^*C).D\).

\end{proposition}
\begin{proof}
The isomorphism away from \(P\) gives the restriction sequence; the kernel is
generated by \(E\), and \(E^2=-1\) proves independence.  Represent divisors on
\(X\) by transverse curves avoiding \(P\), whose strict transforms compute the
first three intersection identities.  The last follows by the projection
formula.
\end{proof}

\begin{proposition}[Canonical divisor under a blow-up]
\label{ha-v3-3}\dcref{V.3.3}\source{hartshorne-algebraic-geometry:V.3.3}
\[
 K_{\widetilde X}=\pi^*K_X+E,\qquad K_{\widetilde X}^2=K_X^2-1.
\]
\uses{ha-v3-2}
\end{proposition}
\begin{proof}
Write \(K_{\widetilde X}=\pi^*K_X+aE\).  Adjunction on
\(E\simeq\mathbb P^1\) gives \(-2=(K_{\widetilde X}+E).E=-a-1\), so \(a=1\).
The square follows from Proposition~\ref{ha-v3-2}.
\end{proof}

\begin{proposition}[Cohomology under a blow-up]
\label{ha-v3-4}\dcref{V.3.4}\source{hartshorne-algebraic-geometry:V.3.4}
\[
 \pi_*\mathcal O_{\widetilde X}=\mathcal O_X,qquad
 R^i\pi_*\mathcal O_{\widetilde X}=0\ (i>0),
\]
so \(H^i(X,\mathcal O_X)\simeq H^i(\widetilde X,\mathcal O_{\widetilde X})\).

\end{proposition}
\begin{proof}
The first equality holds because the blow-up is an isomorphism off a codimension
two point and both schemes are normal.  For the higher direct images, formal
functions reduces to the thickenings \(E_n\) of \(E\).  Their successive
quotients are \(\mathcal O_E(n)\), whose \(H^1\) vanishes for the relevant
nonnegative twists; induction gives \(H^i(E_n,\mathcal O_{E_n})=0\) for
\(i>0\), and formal functions gives the assertion.
\end{proof}

\begin{corollary}[Arithmetic genus]
\label{ha-v3-5}\dcref{V.3.5}\source{hartshorne-algebraic-geometry:V.3.5}
\(p_a(X)=p_a(\widetilde X)\); in particular \(q\) and \(p_g\) are unchanged.
\uses{ha-v3-4}
\end{corollary}
\begin{proof}
Take Euler characteristics in Proposition~\ref{ha-v3-4}.
\end{proof}

\begin{proposition}[Total and strict transforms]
\label{ha-v3-6}\dcref{V.3.6}\source{hartshorne-algebraic-geometry:V.3.6}
If an effective curve \(C\) has multiplicity \(r\) at \(P\), then
\[
 \pi^*C=\widetilde C+rE.
\]

\end{proposition}
\begin{proof}
In the two blow-up charts, divide the local equation of \(C\) by the highest
common power of the exceptional coordinate.  The remaining equation is the
strict transform, and its factor is the initial homogeneous form of degree
\(r\), proving the divisor identity.
\end{proof}

\begin{corollary}[Genus drop on blowing up]
\label{ha-v3-7}\dcref{V.3.7}\source{hartshorne-algebraic-geometry:V.3.7}
With the notation of Proposition~\ref{ha-v3-6},
\[
 \widetilde C.E=r,qquad p_a(\widetilde C)=p_a(C)-\frac12r(r-1).
\]
\uses{ha-v3-6,ha-v3-2}
\end{corollary}
\begin{proof}
The first equality follows from Proposition~\ref{ha-v3-2}.  Substitute
\(\widetilde C=\pi^*C-rE\) and \(K_{\widetilde X}=\pi^*K_X+E\) in adjunction;
the cross terms give exactly \(r(r-1)/2\).
\end{proof}

\begin{proposition}[Resolution of an irreducible curve]
\label{ha-v3-8}\dcref{V.3.8}\source{hartshorne-algebraic-geometry:V.3.8}
Repeatedly blowing up singular points of an irreducible curve produces, after
finitely many steps, a nonsingular strict transform.
\uses{ha-v3-7}
\end{proposition}
\begin{proof}
At each singular point the multiplicity is at least two, so Corollary~\ref{ha-v3-7}
strictly decreases the nonnegative integer arithmetic genus.  The process must
therefore terminate, and termination means that no singular point remains.
\end{proof}

\begin{theorem}[Embedded resolution of curves]
\label{ha-v3-9}\dcref{V.3.9}\source{hartshorne-algebraic-geometry:V.3.9}
For every curve \(Y\subset X\), a finite sequence of monoidal transformations
makes the total inverse image a divisor with normal crossings: every component
is nonsingular, intersections are transverse, and no two components meet in more
than one local branch.
\uses{ha-v3-8}
\end{theorem}
\begin{proof}
Resolve each irreducible component using Proposition~\ref{ha-v3-8}.  If a
remaining singularity has multiplicity \(r\ge3\), blowing it up changes the
arithmetic genus of the total transform by
\(-\frac12(r-1)(r-2)\), so only finitely many such steps occur.  A double point
with distinct tangent directions is already a node; otherwise its blow-up
produces either a separated transverse pair or another configuration of higher
contact, and the next blow-up lowers the same genus measure.  Thus the process
terminates with normal crossings.
\end{proof}

\begin{definition}[Infinitely near points and strict transform]
\source{hartshorne-algebraic-geometry:V.3}
Points on successive blow-ups are called infinitely near points.  A point $Q'$
is near $P$ when it lies on the strict transform of a curve through $P$.
For an effective curve of multiplicity $r$ at $P$, the strict transform is
defined by $\widetilde C=\pi^*C-rE$.
\end{definition}

\section{The Cubic Surface in \(\mathbb P^3\)}
\label{ha-v4}\source{hartshorne-algebraic-geometry:V.4}

\begin{definition}[Assigned base points]
\source{hartshorne-algebraic-geometry:V.4}
An assigned base point of multiplicity $m$ imposes the corresponding vanishing
conditions on a plane linear system.  Points may be infinitely near, in which
case the condition includes the prescribed tangent direction.  Blowing up all
assigned points replaces $|D|$ by
$|\pi^*D-\sum m_iE_i|$; base points remaining after this operation are called
unassigned base points.
\end{definition}

\begin{proposition}[Conics through assigned points]
\label{ha-v4-1}\dcref{V.4.1}\source{hartshorne-algebraic-geometry:V.4.1}
Conics through at most four assigned points in \(\mathbb P^2\), with no three
collinear, have no unassigned base points, including an assigned point infinitely
near another.

\end{proposition}
\begin{proof}
It suffices to consider four points.  The reducible conics
\(L_{12}+L_{34}\) and \(L_{13}+L_{24}\) lie in the system and their common
intersection is exactly the assigned scheme.  In the infinitely near case choose
the corresponding line with the prescribed tangent direction; the same
intersection calculation proves there is no further base point.
\end{proof}

\begin{corollary}[Dimension of conic systems]
\label{ha-v4-2}\dcref{V.4.2}\source{hartshorne-algebraic-geometry:V.4.2}
Under the hypotheses of Proposition~\ref{ha-v4-1}, the system of conics has
dimension \(5-r\) for \(r\le4\); for five points there is a unique irreducible
conic.
\uses{ha-v4-1}
\end{corollary}
\begin{proof}
Each assigned point imposes one independent linear condition on the six-dimensional
space of conics, by Proposition~\ref{ha-v4-1}.  For five points the remaining
one-dimensional vector space cannot contain a reducible conic without forcing
three assigned points onto one line.
\end{proof}

\begin{proposition}[Cubics through assigned points]
\label{ha-v4-3}\dcref{V.4.3}\source{hartshorne-algebraic-geometry:V.4.3}
Cubics through at most seven assigned points, with no four collinear and no seven
on a conic, have no unassigned base points, including infinitely near points.

\end{proposition}
\begin{proof}
For seven ordinary points and a point \(Q\) outside them, choose a line through
two assigned points and a conic through the other five avoiding \(Q\), unless
the incidence forces a line through three points or a conic through all seven;
the hypotheses exclude precisely these exceptions.  Varying the choice gives a
cubic through the assigned points but not \(Q\).  If an assigned point is
infinitely near, impose the tangent condition in the same line--conic
construction; the resulting cubic can be chosen with the required first jet.
\end{proof}

\begin{corollary}[Dimensions of cubic systems]
\label{ha-v4-4}\dcref{V.4.4}\source{hartshorne-algebraic-geometry:V.4.4}
The system of cubics through \(r\le7\) points in the hypotheses of
Proposition~\ref{ha-v4-3} has dimension \(9-r\).  For \(r=8\), it has dimension
one and its general member is irreducible.
\uses{ha-v4-3}
\end{corollary}
\begin{proof}
There are ten cubic monomials, and Proposition~\ref{ha-v4-3} gives independent
conditions.  For eight points the residual pencil has no fixed component,
otherwise a line or conic would contain too many assigned points; Bertini then
gives irreducibility of the general member.
\end{proof}

\begin{corollary}[The ninth base point]
\label{ha-v4-5}\dcref{V.4.5}\source{hartshorne-algebraic-geometry:V.4.5}
Eight points in general position determine a ninth point, possibly infinitely
near, common to every cubic through the eight.

\end{corollary}
\begin{proof}
Two general cubics through the eight have intersection number nine.  They have
no common component and already meet at the eight assigned points, counted with
their assigned tangent conditions, so Bezout leaves one residual point.  Varying
the two cubics shows the residual point is fixed.
\end{proof}

\begin{theorem}[Very ampleness of cubics on the blow-up]
\label{ha-v4-6}\dcref{V.4.6}\source{hartshorne-algebraic-geometry:V.4.6}
If assigned points satisfy no three collinear and no six on a conic, the cubics
through them induce a very ample linear system on their blow-up.
\uses{ha-v4-3}
\end{theorem}
\begin{proof}
After subtracting an exceptional point, the systems needed to separate any two
points or a tangent vector are cubic systems with at most seven assigned points.
Proposition~\ref{ha-v4-3} supplies a member avoiding any prescribed residual
point, and the same argument with a first-order infinitely near point separates
tangent vectors.  Thus the system separates points and tangent vectors.
\end{proof}

\begin{corollary}[Del Pezzo embeddings]
\label{ha-v4-7}\dcref{V.4.7}\source{hartshorne-algebraic-geometry:V.4.7}
Blowing up \(r=0,\ldots,6\) general points of \(\mathbb P^2\) embeds the result
in \(\mathbb P^{9-r}\) with degree \(9-r\), and
\(\omega_{X'}\simeq\mathcal O_{X'}(-1)\).  For \(r=6\), the image is a
nonsingular cubic surface in \(\mathbb P^3\).
\uses{ha-v4-6,ha-v3-2}
\end{corollary}
\begin{proof}
The system is \(3l-\sum E_i=-K_{X'}\).  Theorem~\ref{ha-v4-6} gives the
embedding, while Proposition~\ref{ha-v3-2} computes its square and the
dimension of its sections as \(9-r\).  The canonical blow-up formula gives
\(K_{X'}=-h\), hence the asserted canonical sheaf and degree.
\end{proof}

\begin{proposition}[Picard calculations on a cubic surface]
\label{ha-v4-8}\dcref{V.4.8}\source{hartshorne-algebraic-geometry:V.4.8}
For \(X=\operatorname{Bl}_{P_1,\ldots,P_6}\mathbb P^2\), let \(l\) be the pullback
of a line and \(e_i\) the exceptional classes.  Then
\[
 l^2=1,\ l.e_i=0,\ e_i.e_j=-\delta_{ij},\quad
 h=3l-\sum e_i,\quad K_X=-h,
\]
and for \(D\sim al-\sum b_ie_i\),
\[
 \deg_hD=3a-\sum b_i,\quad D^2=a^2-\sum b_i^2,quad
 p_a(D)=1+\tfrac12(D^2-\deg_hD).
\]
\uses{ha-v3-2,ha-v3-3}
\end{proposition}
\begin{proof}
The intersection table and canonical formula are Propositions~\ref{ha-v3-2}
and~\ref{ha-v3-3}.  The degree is \(D.h\), the square follows by expansion,
and adjunction gives the arithmetic-genus formula.
\end{proof}

\begin{theorem}[The twenty-seven lines]
\label{ha-v4-9}\dcref{V.4.9}\source{hartshorne-algebraic-geometry:V.4.9}
The cubic surface contains exactly 27 lines, namely
\[
 e_i\ (1\le i\le6),\qquad l-e_i-e_j\ (i<j),\qquad
 2l-\sum_{j\ne i}e_j\ (1\le i\le6).
\]
They are all the curves of self-intersection \(-1\).

\end{theorem}
\begin{proof}
A line has \(D.h=1\) and, by adjunction, \(D^2=-1\).  Writing
\(D=al-\sum b_ie_i\), these equations are
\(3a-\sum b_i=1\) and \(a^2-\sum b_i^2=-1\).  Schwarz's inequality gives
\(3a^2-6a-5\le0\), hence \(a=0,1,2\).  Solving the integer equations gives
exactly the six exceptional classes, the fifteen pair-line classes, and the
six conic classes listed above.  Each is effective and represented by the
corresponding plane curve, so the list is complete.
\end{proof}

\begin{proposition}[Six skew lines]
\label{ha-v4-10}\dcref{V.4.10}\source{hartshorne-algebraic-geometry:V.4.10}
Any six mutually disjoint lines on a cubic surface are the exceptional curves
of a morphism \(X\to\mathbb P^2\); thus they define a second blow-up presentation.

\end{proposition}
\begin{proof}
For a chosen six-tuple, use the quadratic transformation of \(\mathbb P^2\) to
replace a selected configuration \((e_i,l-e_i-e_j,2l-\sum_{j\ne i}e_j)\) by
exceptional classes.  The intersection table shows that the remaining six
classes are again disjoint and their images are six points with no three
collinear and no six on a conic.  Induction over the possible incidences gives
the required blow-down map.
\end{proof}

\begin{theorem}[Ampleness on a cubic surface]
\label{ha-v4-11}\dcref{V.4.11}\source{hartshorne-algebraic-geometry:V.4.11}
For a divisor \(D\) on a nonsingular cubic surface, the following are equivalent:
\(D\) is very ample; \(D\) is ample; \(D^2>0\) and \(D.L>0\) for every line
\(L\); and \(D.L>0\) for every line \(L\).
\uses{ha-v4-10,ha-v4-12}
\end{theorem}
\begin{proof}
The implications from very ample to ample and to positivity are immediate;
the converse positivity implies \(D^2>0\) by the intersection description and
Nakai--Moishezon.  Choose six skew lines minimizing their intersections with
\(D\), use Proposition~\ref{ha-v4-10} to relabel them as exceptional classes,
and write \(D=al-\sum b_ie_i\) with ordered positive \(b_i\).  The line
inequalities give \(a\ge b_1+b_2+b_5\).  Lemma~\ref{ha-v4-12} then makes \(D\)
very ample.
\end{proof}

\begin{lemma}[A sufficient very-ampleness decomposition]
\label{ha-v4-12}\dcref{V.4.12}\source{hartshorne-algebraic-geometry:V.4.12}
If \(D\sim al-\sum b_ie_i\), \(b_1\ge\cdots\ge b_6>0\), and
\(a\ge b_1+b_2+b_5\), then \(D\) is very ample.
\uses{ha-v4-6}
\end{lemma}
\begin{proof}
Subtract successively nonnegative multiples of
\[
 l,\quad l-e_1,\quad 2l-e_1-e_2,\quad\ldots,\quad3l-e_1-\cdots-e_6.
\]
The hypotheses ensure nonnegative coefficients.  The first systems are
base-point-free and the last is the anticanonical very ample system from
Theorem~\ref{ha-v4-6}; tensor products and sums of these systems separate
points and tangent vectors, so their nonnegative combination is very ample.
\end{proof}

\begin{corollary}[Numerical criterion on a cubic surface]
\label{ha-v4-13}\dcref{V.4.13}\source{hartshorne-algebraic-geometry:V.4.13}
For \(D\sim al-\sum b_ie_i\), ampleness and very ampleness are equivalent to
\[
 b_i>0,\qquad a>b_i+b_j\ (i\ne j),\qquad
 2a>\sum_{i\ne j}b_i\quad\text{for every }i.
\]
Every such class contains an irreducible nonsingular member.
\uses{ha-v4-11}
\end{corollary}
\begin{proof}
Intersect with the 27 line classes to obtain exactly the displayed inequalities;
Theorem~\ref{ha-v4-11} gives equivalence with ampleness.  A general member is
irreducible and nonsingular by Bertini, since the class is very ample.
\end{proof}

\section{Birational Transformations}
\label{ha-v5}\source{hartshorne-algebraic-geometry:V.5}

\begin{definition}[Birational-map data]
\source{hartshorne-algebraic-geometry:V.5}
For a birational map \(T:X\dashrightarrow Y\), its maximal domain is the
largest open set on which it is a morphism; points outside it are fundamental
points.  Its graph $\Gamma$ is the closure of the graph over this domain.  The
total transform of a subvariety is obtained from the image under the graph
projection, with strict transform away from exceptional components.
\end{definition}

\begin{lemma}[Fundamental points]
\label{ha-v5-1}\dcref{V.5.1}\source{hartshorne-algebraic-geometry:V.5.1}
For a birational map from a normal projective variety, the fundamental locus is
closed and has codimension at least two.

\end{lemma}
\begin{proof}
At a codimension-one point the local ring is a DVR.  The induced inclusion of
function fields defines a valuation on the target field, and the valuative
criterion extends the rational map across that point.  Thus indeterminacy cannot
contain a divisor; closedness follows from the graph projection.
\end{proof}

\begin{theorem}[Zariski's main theorem for a fundamental point]
\label{ha-v5-2}\dcref{V.5.2}\source{hartshorne-algebraic-geometry:V.5.2}
For a birational map of normal projective surfaces, the total transform of each
fundamental point is connected and has dimension at least one.

\end{theorem}
\begin{proof}
Let \(\Gamma\) be the graph and project it to the source.  The fiber over the
fundamental point is connected by the connected-fibers theorem for a proper
birational morphism.  If it were zero-dimensional, the induced projective
birational morphism to a normal affine neighborhood would be finite; a finite
birational morphism to a normal scheme is an isomorphism, contradicting that the
point is fundamental.  Hence the fiber has dimension at least one.
\end{proof}

\begin{proposition}[Factorization through a point blow-up]
\label{ha-v5-3}\dcref{V.5.3}\source{hartshorne-algebraic-geometry:V.5.3}
If \(f:X'\to X\) is a birational morphism of nonsingular projective surfaces,
and \(P\) is a fundamental point of \(f^{-1}\), then \(f\) factors through the
blow-up \(\pi:\widetilde X\to X\) at \(P\): there is \(f':X'\to\widetilde X\)
with \(f=\pi f'\).
\uses{ha-v5-2}
\end{proposition}
\begin{proof}
Set \(T=\pi^{-1}f\).  By Theorem~\ref{ha-v5-2}, the fiber of \(f\) over \(P\)
contains a curve.  In regular local coordinates at a point of that curve, the
pullbacks of the two parameters generate an ideal whose common vanishing is the
fiber; comparing their orders shows that the ideal \(I_P\mathcal O_{X'}\) is
invertible.  The universal property of the blow-up therefore extends \(T\) to
a morphism \(f'\), proving the factorization.
\end{proof}

\begin{corollary}[Birational morphisms]
\label{ha-v5-4}\dcref{V.5.4}\source{hartshorne-algebraic-geometry:V.5.4}
If \(n(f)\) is the number of irreducible curves contracted by a birational
morphism \(f:X'\to X\) of nonsingular projective surfaces, then \(f\) is a
composition of exactly \(n(f)\) monoidal transformations.
\uses{ha-v5-3}
\end{corollary}
\begin{proof}
Choose a contracted curve over a fundamental point of the inverse and apply
Proposition~\ref{ha-v5-3}; the factorization removes one contracted component.
Induction on \(n(f)\) terminates, and each blow-up adds exactly one exceptional
curve.
\end{proof}

\begin{theorem}[Factorization of birational transformations]
\label{ha-v5-5}\dcref{V.5.5}\source{hartshorne-algebraic-geometry:V.5.5}
Every birational transformation between nonsingular projective surfaces is a
composition of monoidal transformations and their inverses.
\uses{ha-v3-7,ha-v5-4}
\end{theorem}
\begin{proof}
Choose a very ample divisor \(H'\) on the target and a nonsingular
\(C'\in|2H'|\) avoiding the fundamental points.  Pull it back to the source;
the difference \(m=p_a(C)-p_a(C')\) is nonnegative.  Blow up singular and
infinitely near points of the pullback curve.  Corollary~\ref{ha-v3-7} strictly
decreases \(m\), so after finitely many steps the rational map becomes a
morphism.  Apply Corollary~\ref{ha-v5-4} to that morphism and to the analogous
map for the inverse.
\end{proof}

\begin{corollary}[Birational invariance of surface arithmetic genus]
\label{ha-v5-6}\dcref{V.5.6}\source{hartshorne-algebraic-geometry:V.5.6}
The arithmetic genus of a nonsingular projective surface is a birational
invariant.
\uses{ha-v5-5,ha-v3-5}
\end{corollary}
\begin{proof}
Theorem~\ref{ha-v5-5} reduces a birational map to blow-ups and blow-downs, and
Corollary~\ref{ha-v3-5} preserves arithmetic genus at each step.
\end{proof}

\begin{theorem}[Castelnuovo contraction]
\label{ha-v5-7}\dcref{V.5.7}\source{hartshorne-algebraic-geometry:V.5.7}
If \(Y\simeq\mathbb P^1\) on a nonsingular projective surface and \(Y^2=-1\),
there is a morphism contracting \(Y\) to a nonsingular point, and this morphism
is inverse to the blow-up of that point.

\end{theorem}
\begin{proof}
Choose very ample \(H\) with \(H^1(H)=0\), put \(k=H.Y\), and use
\(0\to\mathcal O(H+(i-1)Y)\to\mathcal O(H+iY)\to\mathcal O_Y(H+iY)\to0\)
to prove inductively \(H^1(H+iY)=0\) for \(0\le i\le k\).  Consequently
\(\mathcal O(H+kY)\) is globally generated.  Its morphism is an isomorphism
off \(Y\) and contracts \(Y\).  Normalize the image.  Formal functions and
the thickenings of \(Y\) identify the completed local ring at the image with
\(\varprojlim k[x,y]/(x,y)^{n+1}=k[[x,y]]\), so the point is nonsingular.
Factorization identifies the morphism with the blow-down; the same completed
ring calculation shows the normalization map is already an isomorphism.
\end{proof}

\begin{theorem}[Relatively minimal models]
\label{ha-v5-8}\dcref{V.5.8}\source{hartshorne-algebraic-geometry:V.5.8}
Every nonsingular projective surface admits a birational morphism to a relatively
minimal surface, i.e. one with no exceptional curve of the first kind.
\uses{ha-v5-7}
\end{theorem}
\begin{proof}
Repeatedly contract every \((-1)\)-curve by Theorem~\ref{ha-v5-7}.  The
cohomology classes of contracted curves are linearly independent in
\(H^1(X,\Omega_X)\), since their intersection matrix is negative diagonal after
successive pullback.  Their number is therefore bounded by the finite dimension
of this space, so the process terminates at a relatively minimal surface.
\end{proof}

\section{Classification of Surfaces}
\label{ha-v6}\source{hartshorne-algebraic-geometry:V.6}

\begin{definition}[Canonical ring and Kodaira dimension]
\source{hartshorne-algebraic-geometry:V.6}
For a nonsingular projective surface define the canonical ring
\(R(X,K)=\bigoplus_{n\ge0}H^0(X,\mathcal O_X(nK_X))\).  Its Kodaira dimension
\(\kappa(X)\) is \(-1\) when all positive pluricanonical systems are empty;
otherwise it is the dimension of the image of \(|nK_X|\) for sufficiently
divisible \(n\), and lies in \(\{0,1,2\}\).
\end{definition}

\begin{theorem}[The case \(\kappa=-1\)]
\label{ha-v6-1}\dcref{V.6.1}\source{hartshorne-algebraic-geometry:V.6.1}
\[
 \kappa(X)=-1\quad\Longleftrightarrow\quad |12K_X|=\varnothing
 \quad\Longleftrightarrow\quad X\text{ is rational or ruled}.
\]

\end{theorem}
\begin{proof}
The source states this classification theorem but does not supply its proof in
the chapter body.
\end{proof}

\begin{theorem}[Castelnuovo rationality criterion]
\label{ha-v6-2}\dcref{V.6.2}\source{hartshorne-algebraic-geometry:V.6.2}
For a nonsingular projective surface, \(X\) is rational if and only if
\(p_a(X)=0\) and \(P_2=0\), where
\(P_2=h^0(X,\mathcal O_X(2K_X))\).

\end{theorem}
\begin{proof}
The source cites Castelnuovo's criterion and does not supply its proof here.
\end{proof}

\begin{theorem}[The case \(\kappa=0\)]
\label{ha-v6-3}\dcref{V.6.3}\source{hartshorne-algebraic-geometry:V.6.3}
If \(\operatorname{char}k\ne2,3\), then
\(\kappa(X)=0\) iff \(|12K_X|\) consists of one divisor (equivalently
\(P_{12}=1\)).  The minimal surfaces with \(\kappa=0\) are K3 surfaces,
Enriques surfaces, abelian surfaces, and hyperelliptic (bielliptic) surfaces,
distinguished by their canonical torsion and invariants \(p_g,q\).

\end{theorem}
\begin{proof}
The source states the classification and does not supply its proof in this
chapter body.
\end{proof}

\begin{theorem}[The case \(\kappa=1\)]
\label{ha-v6-4}\dcref{V.6.4}\source{hartshorne-algebraic-geometry:V.6.4}
\(\kappa(X)=1\) if and only if \(X\) is an elliptic surface: after a birational
modification it admits a fibration to a curve whose general fiber is a smooth
curve of genus one.

\end{theorem}
\begin{proof}
The source states this characterization and does not supply its proof here.
\end{proof}

\begin{theorem}[The case \(\kappa=2\)]
\label{ha-v6-5}\dcref{V.6.5}\source{hartshorne-algebraic-geometry:V.6.5}
\(\kappa(X)=2\) if and only if some pluricanonical system \(|nK_X|\) defines a
birational map onto its image.  Such surfaces are precisely the surfaces of
general type.

\end{theorem}
\begin{proof}
The source states this characterization and does not supply its proof here.
\end{proof}
